\documentclass[10pt,a4paper]{amsart}
\usepackage[margin=19mm]{geometry}
\usepackage{amsmath,amssymb,mathtools}
\usepackage[scr=rsfs,cal=euler]{mathalfa}
\usepackage{xfrac}
\usepackage[unicode,hidelinks,bookmarks=false]{hyperref}
\newcommand{\ie}{i.e.\ }
\newcommand{\cf}{cf.\ }
\newcommand{\resp}{respectively\ }
\newcommand{\Subsection}{\subsection}
\providecommand{\nobreakdash}{\nobreak\hbox{-}}
\setlength{\emergencystretch}{2em}
\multlinegap0pt
\usepackage[shortlabels]{enumitem}
\usepackage{amssymb}
\usepackage{mathtools}
\usepackage{tikz}
\usepackage{xcolor}
\newtheorem{theorem}{Theorem}
\newtheorem{fact}{Fact}
\newtheorem{lemma}{Lemma}
\newcommand{\PP}{\mathbb{P}}
\title{Grid-free linear hypergraphs via Cayley-Bacharach}
\author{Cosmin Pohoata}
\date{Source: arXiv:2602.14716v1 (CC BY 4.0)}
\begin{document}
\maketitle

\noindent\emph{Source and licence.} Grid-free linear hypergraphs via Cayley-Bacharach. Version arXiv:2602.14716v1 (CC BY 4.0), distributed by arXiv under the Creative Commons Attribution 4.0 International licence, http://creativecommons.org/licenses/by/4.0/, as recorded in the arXiv licence metadata for this exact version and shown on the abstract page https://arxiv.org/abs/2602.14716v1. Full licence notice: \texttt{licenses/arxiv-2602.14716-CC-BY-4.0.txt}.\par\smallskip

\noindent\textit{Editorial scope.} Counted unit: the article's theorem thm:main (source0.tex lines 108--123). The article's complete original proof of it is delivered with it (source0.tex lines 314--363, one \texttt{proof} environment of 50 lines, which the article places away from the statement). No mathematics is written, repaired or re-derived: the statement and the proof are the article's own lines, in the article's own notation, and no retained line is substituted or re-flowed.  Retained uncounted as the source-local context the proof needs: lines 231--233; lines 241--243; lines 303--308.  Nothing was omitted from any retained span, so the package carries no omission record.  No retained line contains a citation, so the wrapper carries no bibliography and no bibliography entry.  No retained line contains a figure environment, an image-inclusion command or drawing code, so no image is delivered and the package declares no asset dependency.  Editorial formatting only, in addition: an AMS \texttt{amsart} preamble that restates the article's own declarations and macros used by the retained lines, namely the environments \texttt{theorem}, \texttt{fact}, \texttt{lemma} on separate counters, together with the standard packages those lines need.  The retained environments print in this document as Theorem 1, Fact 1, Lemma 1 in order of appearance, whereas the article prints the same items with the numbering of its own full document.  The complete native source of the article as distributed in the arXiv e-print for this exact version is bundled unchanged as \texttt{sources/194737/source0.tex}.
\begin{fact}\label{lem:disjoint}
$A\cap B=\varnothing$.
\end{fact}

\begin{fact}\label{lem:nothree}
No affine line meets $B$ in three distinct points.
\end{fact}

\begin{lemma}[Cayley-Bacharach]\label{cor:planeCB}
Let $D_1,D_2\subset \PP^2$ be plane curves of degrees $r$ and $r$ with no common component.
Assume $D_1\cap D_2$ consists of exactly $r^2$ distinct points $X=\{x_1,\dots,x_{r^2}\}$.
Then any homogeneous polynomial $f$ of degree $\le 2r-3$ that vanishes on $r^2-1$ points of $X$
must vanish on all $r^2$ points of $X$.
\end{lemma}

\begin{theorem}\label{thm:main}
For every $r \geq 3$ and every sufficiently large odd prime power $q$, there exists an $r$-uniform linear
hypergraph $H_{r,q}$ with no $r\times r$ grid, with
\[
|V(H_{r,q})| = rq,
\qquad
|E(H_{r,q})| = \frac{q^2+2q-1}{2}.
\]
Equivalently, for $n=rq$,
\[
|E(H_{r,q})|
=
\left(\frac{1}{2r^2}+o(1)\right)n^2
\qquad (q\to\infty,\ r\text{ fixed}).
\]
\end{theorem}

\begin{proof}[Proof of Theorem \ref{thm:main}]
Assume for contradiction that $H_{r,q}$ contains a subhypergraph isomorphic to $G_{r\times r}$.
Let $R_1,\dots,R_r$ be the row edges and $C_1,\dots,C_r$ the column edges.

By construction, every edge of $H_{r,q}$ is supported on a unique nonhorizontal affine line.
Thus there exist affine lines $\ell_1,\dots,\ell_r$ and $m_1,\dots,m_r$ such that
\[
R_i=e(\ell_i)\quad(1\le i\le r),
\qquad
C_j=e(m_j)\quad(1\le j\le r).
\]
No $\ell_i$ equals any $m_j$: otherwise $R_i=C_j$ would have $r$ common vertices, contradicting $|R_i\cap C_j|=1$.

Because $|R_i\cap C_j|=1$, the lines $\ell_i$ and $m_j$ intersect in the affine plane (so they are not parallel),
and their intersection point is exactly the grid vertex $R_i\cap C_j$.
Since the grid has $r^2$ distinct vertices, the $r^2$ points $\ell_i\cap m_j$ are all distinct. Let $\overline{\ell_i},\overline{m_j}\subset \PP^2$ be projective closures and define degree-$r$ curves
\[
D_1 := \overline{\ell_1}\cup\cdots\cup \overline{\ell_r},
\qquad
D_2 := \overline{m_1}\cup\cdots\cup \overline{m_r}.
\]
Then $D_1\cap D_2$ consists of exactly the $r^2$ distinct points $\overline{\ell_i}\cap \overline{m_j}$
(all affine), hence Lemma ~\ref{cor:planeCB} applies to $D_1,D_2$.

Every edge of $H_{r,q}$ contains exactly one vertex in $B$.
The grid has $2r$ edges, hence $2r$ incidences of the form ``(edge, its $B$-vertex)''.
In the grid, every vertex lies in exactly two edges (one row, one column), so each $B$-vertex is counted twice.
Therefore the grid contains exactly $(2r)/2=r$ distinct vertices on $B$.
Label these points $b_1,\dots,b_r\in B$.
All remaining $r^2-r=r(r-1)$ grid vertices lie in $A$. Define the curve
\[
D \;:=\; A\ \cup\ \overline{b_1b_2}\ \cup\ \overline{b_2b_3}\ \cup\ \cdots\ \cup\ \overline{b_{r-2}b_{r-1}},
\]
a union of $(r-1)$ lines (those forming $A$) and $(r-2)$ additional lines, hence its degree satisfies $\deg(D)=2r-3$.
Let $f$ be a homogeneous polynomial of degree $2r-3$ defining $D$.

By construction:
\begin{itemize}[leftmargin=2em]
\item $D$ contains every grid vertex lying in $A$, since $A\subset D$.
\item $D$ contains $b_1,\dots,b_{r-1}$ (each is an endpoint of one of the lines $\overline{b_ib_{i+1}}$).
\item $b_r\notin D$: indeed $b_r\notin A$ by Fact~\ref{lem:disjoint}, and if $b_r$ lay on
$\overline{b_ib_{i+1}}$ then that line would meet the conic $B$ in the three distinct points
$b_i,b_{i+1},b_r$, contradicting Fact~\ref{lem:nothree}.
\end{itemize}
Hence $f$ vanishes on all grid vertices except $b_r$. In other words, if $X=D_1\cap D_2$ is the set of $r^2$ grid vertices, then $f$ vanishes on $r^2-1$ points of $X$
but not on the last point $b_r$.

Since $\deg f = 2r-3$, Lemma~\ref{cor:planeCB} forces $f$ to vanish on \emph{all} points of $X$,
in particular on $b_r$ as well, contradiction. Therefore $H_{r,q}$ contains no $r\times r$ grid.
\end{proof}

\end{document}
